\uuid{pc5l}
\titre{Base orthonormée}
\auteur{bodin ; hanani ; mzari}
\organisation{exo7}
\licence{CC BY-NC-SA 4.0}
\datecreate{2020-01-20}
\niveau{L1}
\module{Géométrie}
\chapitre{Géométrie du plan}
\sousChapitre{Distance, norme, produit scalaire}
\difficulte{2}
\theme{base orthonormée}
\exo7qcmid{lille-1.geometrie-plane.012}

\contenu{
\begin{question}
\qclassification{140.01,140.02}
\qtags{motcle=base orthonormée}
Dans le plan muni d'une base orthonormée $(\vec{i},\vec{j})$, on considère les vecteurs $\displaystyle \vec{u}=\left(\frac{1}{\sqrt{2}},a\right)$ et $\displaystyle \vec{v}=\left(a,-\frac{1}{\sqrt{2}}\right)$. Comment choisir le réel $a$ pour que $(\vec{u},\vec{v})$ soit une base orthonormée ?
\begin{answers}
    \good{$\displaystyle a=\frac{1}{\sqrt{2}}$}
    \good{$\displaystyle a=-\frac{1}{\sqrt{2}}$}
    \bad{$\displaystyle a=\sqrt{2}$}
    \bad{$\displaystyle a=-\sqrt{2}$}
\end{answers}
\begin{explanations}
Pour tout $a\in \Rr$, les vecteurs $\vec{u}$ et $\vec{v}$ sont orthogonaux. Ensuite, $\|\vec{u}\|=\|\vec{v}\|=1$ implique $\displaystyle a=\frac{\pm 1}{\sqrt{2}}$.
\end{explanations}
\end{question}
}
